% !PRINTED 627
\weberpage{627}

%<<CONT
und wenn man hierin f\"ur $\lambda$ und $\mu$ rationale Gr\"ossen setzt (aus
einem beliebigen Rationalit\"atsbereich, z.\,B. auch aus dem K\"orper
der rationalen Functionen von $\lambda$ und $\mu$), so ist die Gleichung
5$^{\text{ten}}$ Grades
\weq{24}{x^5+\alpha x+\beta=0}
algebraisch l\"osbar, da ja der Fall, dass diese Gleichung reducibel
 wird, auch auf algebraisch l\"osbare Gleichungen f\"uhrt.

Umgekehrt k\"onnen wir sagen, dass keine irreducible Gleichung
 5$^{\text{ten}}$ Grades von der Form (24) algebraisch l\"osbar ist, in
der die Co\"efficienten $\alpha$, $\beta$ nicht in der Form (23) darstellbar sind.

Setzen wir z.\,B. $\lambda=-1$, $\mu=1$, so erhalten wir $64\alpha=-5^4$,
$64\beta=-5^4$. Die Gleichung (24) vereinfacht sich, wenn wir f\"ur
$x$ eine neue Unbekannte $\xi$ einf\"uhren, durch die Gleichung $x\xi=5$.
Wir erhalten so die Gleichung
\[
\xi^5+5\xi^4-5\cdot 64=0
\]
als Beispiel einer \gesperrt{algebraisch l\"osbaren} Gleichung, die \"uberdies
 nach dem vorhin erw\"ahnten Satze \gesperrt{irreducibel} ist.

\begin{center}
\S.\,182.

\gesperrt{Die Gruppe der Resolvente.}
\end{center}

Die S\"atze, die wir im vorigen Paragraphen abgeleitet haben,
gestatten einen merkw\"urdigen Schluss \"uber die Gleichungen
6$^{\text{ten}}$ Grades.

Wir nehmen jetzt wieder die $x_0$, $x_1$, $x_2$, $x_3$, $x_4$ des vorigen
Paragraphen als unabh\"angige Variable an. Dann sind die sechs
Gr\"ossen
\weq{1}{\begin{aligned}
v_1&=(u_1-u_1')^2, & v_2&=(u_2-u_2')^2, & v_3&=(u_3-u_3')^2\\
v_4&=(u_4-u_4')^2, & v_5&=(u_5-u_5')^2, & v_6&=(u_6-u_6')^2
\end{aligned}}
die Wurzeln einer Gleichung 6$^{\text{ten}}$ Grades, $F(v)=0$, deren
Co\"efficienten rational durch die symmetrischen Grundfunctionen
$a$, $b$, $c$, $d$, $e$ der $x$ ausdr\"uckbar sind, und die in dem K\"orper
der rationalen Functionen dieser Gr\"ossen irreducibel ist. Diese
Gleichung ist eine Resolvente der Gleichung 5$^{\text{ten}}$ Grades, $f(x)=0$,
deren Wurzeln die $x$ sind.

Ist $M$ die volle metacyklische Gruppe der $x$, so haben die
zu $M$ conjugirten Gruppen $\pi^{-1}M\pi$ keinen anderen gemeinschaft-
%CONT>>
